\chapter{ANSI serialization and the writer}
\label{ch:ansi}

The diff in the previous chapter tells us \emph{which cells changed
and what they should be}. This chapter is about the last mile: turning
that information into the exact byte stream the terminal expects, and
getting those bytes out to the wire as cheaply as possible.

\section{What the renderer must emit}

For every run of changed cells, the renderer must produce:

\begin{enumerate}[leftmargin=*]
  \item A \textbf{cursor positioning} sequence to move to the start
    of the run.
  \item A \textbf{style change} (SGR) sequence if the new style
    differs from the cursor's current style.
  \item The \textbf{UTF-8 bytes} of the cells' code points.
  \item Optionally, \textbf{hyperlink} OSC-8 wrap if a link is
    associated with the cells.
\end{enumerate}

For a typical animation frame at, say, 1\% changed cells on an
80x24 grid, that is around 20 runs, each maybe 4 cells long, each
requiring a CUP (8-10 bytes), maybe an SGR (15-30 bytes for full
truecolour), and the UTF-8 bytes. Total: about a kilobyte per frame.
That kilobyte must be assembled and written in well under one frame
($\sim$16 ms) to keep up with 60 FPS.

\section{Zero-allocation building blocks}

The old version of \maya{}'s renderer built a \code{std::vector<RenderOp>}
and then serialised it --- which allocated for every entry. The
current version writes ANSI bytes directly into one growing
\code{std::string}. Every per-cell allocation was eliminated.

The lowest layer is \file{maya/include/maya/terminal/ansi.hpp}, a
collection of free functions that append ANSI sequences to a
caller-supplied string:

\begin{lstlisting}
MAYA_FORCEINLINE void write_cursor_up(std::string& out, int n);
MAYA_FORCEINLINE void write_cursor_down(std::string& out, int n);
MAYA_FORCEINLINE void write_move_to(std::string& out, int col, int row);
MAYA_FORCEINLINE void write_erase_lines(std::string& out, int n);
\end{lstlisting}

Each one is a few \code{out += ...} calls and a single
\code{std::to\_chars} for the integer parameter. \code{std::to\_chars}
is \cpp{}17's allocation-free, locale-free integer formatter; it
writes into a stack buffer that we then \code{append} onto \code{out}.

\subsection{A worked function}

\begin{lstlisting}
namespace detail {
    MAYA_FORCEINLINE void append_int(std::string& out, int n) {
        char buf[12];                            // 2^31 fits in 11 chars + sign
        auto [end, _] = std::to_chars(buf, buf + sizeof(buf), n);
        out.append(buf, end);
    }
}

MAYA_FORCEINLINE void write_move_to(std::string& out, int col, int row) {
    out += "\x1b[";
    detail::append_int(out, row);
    out += ';';
    detail::append_int(out, col);
    out += 'H';
}
\end{lstlisting}

No heap allocation. No \code{std::format}. \code{out += "\textbackslash
x1b["} appends a \code{string\_view} which is two characters; the
integers add 1--3 bytes; the totals are usually 8-10 bytes per CUP.
\code{MAYA\_FORCEINLINE} is a portable wrapper for
\code{\_\_attribute\_\_((always\_inline))} (GCC/Clang) and
\code{\_\_forceinline} (MSVC) --- a hint that this should be inlined
even at low optimisation levels.

\subsection{Per-cell positioning, faster}

The diff loop's hot path bypasses even \code{std::to\_chars} for the
common case of row and column numbers under 1000, using a hand-written
ASCII digit splitter:

\begin{lstlisting}
MAYA_FORCEINLINE char* write_uint_pos(char* p, unsigned n) noexcept {
    if (n < 10) {
        *p++ = char('0' + n);
    } else if (n < 100) {
        *p++ = char('0' + n / 10);
        *p++ = char('0' + n % 10);
    } else if (n < 1000) {
        *p++ = char('0' + n / 100);
        *p++ = char('0' + (n / 10) % 10);
        *p++ = char('0' + n % 10);
    } else {
        // fallback: std::to_chars
    }
    return p;
}
\end{lstlisting}

Inlined into the diff loop, this compiles to a tiny chain of
conditional moves and stores --- a few cycles per integer. Premature
optimisation? Not in a loop that runs once per changed cell.

\section{UTF-8 encoding}

For every changed cell we need to encode its 32-bit code point into
1--4 UTF-8 bytes. \maya{}'s encoder is in \code{render/diff.hpp}:

\begin{lstlisting}
MAYA_FORCEINLINE void encode_utf8(char32_t cp, std::string& out) {
    if (cp < 0x80) [[likely]] {                      // ASCII fast path
        out += char(cp);
        return;
    }
    char buf[4];
    int len;
    if (cp < 0x800) {                                // 2-byte
        buf[0] = char(0xC0 | (cp >> 6));
        buf[1] = char(0x80 | (cp & 0x3F));
        len = 2;
    } else if (cp < 0x10000) [[likely]] {            // 3-byte (BMP)
        buf[0] = char(0xE0 |  (cp >> 12));
        buf[1] = char(0x80 | ((cp >> 6) & 0x3F));
        buf[2] = char(0x80 |  (cp        & 0x3F));
        len = 3;
    } else {                                         // 4-byte (supplementary)
        buf[0] = char(0xF0 |  (cp >> 18));
        buf[1] = char(0x80 | ((cp >> 12) & 0x3F));
        buf[2] = char(0x80 | ((cp >>  6) & 0x3F));
        buf[3] = char(0x80 |  (cp        & 0x3F));
        len = 4;
    }
    out.append(buf, len);
}
\end{lstlisting}

The \code{[[likely]]} attribute tells the compiler to lay out the
branches with the ASCII and BMP cases on the hot path; the rare
4-byte case takes the cold branch.

\section{The style cache}

A style change emits an SGR sequence whose exact bytes depend on the
\code{Style} struct: \code{ESC[0;1;3;38;2;100;180;255;48;2;20;25;35m}
for ``reset, bold, italic, fg=(100,180,255), bg=(20,25,35)''. Building
that string takes a dozen \code{to\_chars} calls and several
conditional branches.

Doing this once per changed cell would dominate runtime. \maya{} does
it \emph{once per distinct style, ever} --- when the style is
interned into the \code{StylePool} --- and caches the bytes:

\begin{lstlisting}
StyleId StylePool::intern(Style s) {
    auto [it, inserted] = lookup_.try_emplace(s, table_.size());
    if (inserted) {
        table_.push_back(s);
        cache_sgr_.push_back(build_sgr_string(s));   // do it once
    }
    return it->second;
}
\end{lstlisting}

The diff loop's style change is then a single \code{memcpy}:

\begin{lstlisting}
if (new_style_id != current_style_id) {
    out += pool.sgr(new_style_id);
    current_style_id = new_style_id;
}
\end{lstlisting}

No formatting. No branches. The bytes are already laid out.

\section{The terminal writer}

After the renderer assembles its kilobyte of ANSI, the runtime needs
to put those bytes on the terminal file descriptor. The writer
(\file{maya/include/maya/terminal/writer.hpp}) wraps that:

\begin{lstlisting}
class Writer {
public:
    explicit Writer(int fd) : fd_(fd) {}
    void write(std::string_view bytes);
    void flush();
private:
    int         fd_;
    std::string buffer_;
};
\end{lstlisting}

Three details worth noting:

\begin{itemize}[leftmargin=*]
  \item \textbf{Buffered.} Multiple \code{write()} calls accumulate
    into \code{buffer\_}; one \code{::write(fd, ...)} system call
    sends them all. System calls are expensive; one big write beats
    many small writes.
  \item \textbf{Partial-write loop.} A blocking \code{::write} may
    return having sent fewer bytes than requested. The writer loops
    until everything is out (or an error pins it).
  \item \textbf{Single-frame atomicity.} The whole frame is one
    contiguous string. Even if the terminal is slow, the user never
    sees a half-rendered frame --- the writer hands the whole
    buffer to the kernel in one go.
\end{itemize}

\subsection{Synchronised output mode}

A subtle problem: even with one \code{write}, the terminal can paint
the bytes incrementally, producing a brief flash of the half-applied
frame (a phenomenon called \emph{tearing}). Modern emulators support
``synchronised output mode'' --- a pair of DCS escape sequences that
tell the emulator ``hold off painting until I say I am done.''
\maya{} wraps each frame in these:

\begin{lstlisting}
ESC P = 1 s ESC \    // begin synchronised update
... frame bytes ...
ESC P = 2 s ESC \    // end synchronised update; paint atomically
\end{lstlisting}

Emulators that do not understand the sequence ignore it; those that
do paint the frame as one unit. Result: zero tearing on supporting
terminals, no regression on others.

\section{Putting it all together}

Here is the simplified inner loop of the renderer's \code{emit\_run}:

\begin{lstlisting}
void emit_run(std::string& out, int x, int y,
              const uint64_t* cells, std::size_t count,
              const StylePool& pool,
              StyleId& current_style)
{
    // 1. Cursor positioning (1-based in ANSI; +1 to both).
    out += "\x1b[";
    char buf[16];
    char* p = buf;
    p = write_uint_pos(p, y + 1);
    *p++ = ';';
    p = write_uint_pos(p, x + 1);
    *p++ = 'H';
    out.append(buf, p - buf);

    // 2. Cells.
    for (std::size_t i = 0; i < count; ++i) {
        auto [cp, sty, width] = unpack(cells[i]);
        if (sty != current_style) {
            out += pool.sgr(sty);
            current_style = sty;
        }
        if (width == 0) continue;       // trailing half of a wide cell
        encode_utf8(cp, out);
    }
}
\end{lstlisting}

Notice the discipline: every operation is either a \code{string\_view}
append, a \code{memcpy}, or a few stack-buffer writes followed by an
append. There are no allocations inside the loop.

\section{Cost, measured}

On a modern x86 laptop, with AVX2 and a 200x50 terminal, \maya{}'s
render loop --- including layout, paint, diff, and serialise ---
typically takes \emph{30-80 microseconds} per frame in steady state.
At 60 FPS the per-frame budget is 16,666 microseconds; we are using
under 1\%. The remaining headroom is the terminal emulator's
business.

\begin{recap}
The renderer turns the diff into ANSI bytes via a chain of
allocation-free helpers: stack-buffered integer formatting,
inline-friendly UTF-8 encoding, and a per-style pre-cached SGR
string. The writer batches the frame into one \code{::write()} and
optionally wraps it in synchronised-output mode for tear-free
rendering. The hot path is dominated by \code{memcpy}; the
allocation count is zero.
\end{recap}

\section*{Workshop}

\begin{enumerate}[leftmargin=*]
  \item Read \file{maya/include/maya/terminal/ansi.hpp}. Identify
    every \code{constexpr string\_view} constant and match it to an
    ANSI sequence from Chapter~\ref{ch:terminal}.
  \item In a scratch program, construct a \code{Canvas}, write some
    text into it with \code{draw\_text}, and dump the diff bytes
    against an empty front buffer. Inspect the bytes by hand.
  \item Toggle the synchronised-output sequences off in a debug
    fork of \maya{} and run an animation example. On a fast
    terminal you will see frames tear; on a slow one (try
    \code{ssh localhost} into the same machine) the effect is
    more obvious.
\end{enumerate}
